\scanpage{51}
\mspage{50}
$x$\unsure{''}. i.e. $\nexists\,Y\in M$ \unsure{avec} $x\leq Y<X$. C'est évidemment une condition des plus restrictives sur $M$ ! Passons, et ne nous préoccupons pas de At 8.
% slanted
\margin{``tout élément de $M$ est minimal, ou successeur d'un minimal''. Sous condition des \unsure{chaînes}, cela implique $\dim M\leq1$.}
Considérons \deleted{des} un autre $(M',\leq,\ll\overset{\mathrm{dfn}}{=}\leq,F')$, prenons \deleted{\unsure{tt}} \unsure{pour|pour y} \unsure{faire|fixer} $\unsure{\mathcal F|F|\mathscr F}=\unsure{Pf|P_f}(M)$, $\mathcal F'=Pf(M')$, et intéressons-nous aux
\[
f:F'\longrightarrow F
\]
\deleted{tels} qui soient des morphismes d'ateliers, au sens de la définition p. 47. Par le lemme, ils correspondent aux \deleted{\unsure{hom}} applications croissantes en sens inverse
\diagram{\begin{tikzcd}[ampersand replacement=\&,row sep=small,column sep=2em]
M \arrow[d,"\unsure{\varphi|\psi}"]\\
M' \& {,}
\end{tikzcd}}
sans aucune condition supplémentaire.

Cela montre que la condition que $\varphi$ de (1.55) soit telle que la top. de $\widetilde F$ soit \deleted{\unsure{l'image}} initiale pour $\varphi$, i.e. que (1.54) soit tjrs surjective, est bien substantielle.

\uline{Exemples} 1) Soit $f:L\unsure{-|\to}L'$ application ensembliste, elle induit \deleted{\unsure{Fig}}
\[
\begin{aligned}
f^{\unsure{*|x}}:\operatorname{Fig}(L')&-\operatorname{Fig}(L)\\
\deleted{\unsure{F'|f^{?}}}\deleted{\unsure{A'}}\unsure{\Phi'|\underline{\Phi}'}&\longmapsto\{f^{-1}(A')\mid A'\in\Phi',\, f^{-1}(A')\ne\varnothing\}
\end{aligned}
\]
